\chapter{Graphklassen unter \textit{preferred}-Unabh\"angigkeit}
\label{ch:graphclasses-preferred}

\section{Einleitung}
\label{sec:gcp-einleitung}

Die Kapitel~\ref{ch:graphclasses} und~\ref{ch:graphclasses-empirical} haben die
bedingte Unabh\"angigkeit unter der \textit{stable}-Semantik aus theoretischer
und empirischer Sicht nach Graphklassen aufgeschl\"usselt. Dieses Kapitel
\"ubertr\"agt die Fragestellung auf die \textit{preferred}-Semantik. Wir
verzichten dabei bewusst auf eine empirische Untersuchung und konzentrieren uns
auf konkrete, beweisbare Aussagen.

Gegen\"uber der \textit{stable}-Semantik weist die
\textit{preferred}-Semantik drei Eigenschaften auf, die das Bild verschieben:
\begin{enumerate}
    \item \textbf{Existenz.} Jedes (endliche) AF besitzt mindestens ein
    \textit{preferred}-Labelling, w\"ahrend
    \(\Lambda_{\mathrm{st}}(F)=\emptyset\) m\"oglich ist
    \cite{dung:1995:af}. Der in Kapitel~\ref{ch:graphclasses-empirical}
    pr\"agende Fall \enquote{vakuose Unabh\"angigkeit mangels Labellings}
    entf\"allt damit f\"ur die \textit{preferred}-Semantik vollst\"andig.
    \item \textbf{Maximalit\"at.} Ein \textit{preferred}-Labelling ist ein
    \textit{complete}-Labelling mit \(\subseteq\)-maximaler In-Menge
    (Definition~\ref{def:sem-labelling}). Diese zus\"atzliche Bedingung
    erschwert die Kodierung (Kapitel~\ref{ch:preferred}).
    \item \textbf{Komplexit\"at.} Das Entscheidungsproblem
    \(\textit{Ind}_{\mathrm{pr}}\) ist \(\Pi_3^{\mathrm P}\)-vollst\"andig
    (Theorem~\ref{th:ind-pr-is-pi3p-c}) und damit eine Stufe der polynomiellen
    Hierarchie h\"oher angesiedelt als \(\textit{Ind}_{\mathrm{st}}\)
    (Theorem~\ref{th:ind-st-is-pi2p-c}).
\end{enumerate}

Wir erinnern an die Definition der bedingten Unabh\"angigkeit
(Definition~\ref{def:BU:AFs}), angewandt auf \(\sigma=\mathrm{pr}\):
\(X\bot_{\mathrm{pr}}Y\mid Z\) gilt genau dann, wenn zu je zwei
\textit{preferred}-Labellings \(\lambda_1,\lambda_2\in\Lambda_{\mathrm{pr}}(F)\)
mit \(\lambda_1|_Z=\lambda_2|_Z\) ein
\(\lambda_3\in\Lambda_{\mathrm{pr}}(F)\) mit \(\lambda_3|_X=\lambda_1|_X\),
\(\lambda_3|_Y=\lambda_2|_Y\) und \(\lambda_3|_Z=\lambda_1|_Z\) existiert.

Der Aufbau des Kapitels folgt der Beobachtung, dass sich die
\textit{preferred}-Semantik auf mehreren wichtigen Graphklassen entweder auf
ein \emph{einziges} Labelling reduziert oder mit der \textit{stable}-Semantik
\emph{zusammenf\"allt}. Abschnitt~\ref{sec:gcp-mittel} stellt dazu zwei
allgemeine Beweismittel bereit; die folgenden Abschnitte wenden sie klassenweise
an.

\section{Zwei \"ubergreifende Beweismittel}
\label{sec:gcp-mittel}

\subsection{Das Transferlemma}
\label{subsec:gcp-transfer}

Die bedingte Unabh\"angigkeit h\"angt nach Definition~\ref{def:BU:AFs}
ausschlie\ss lich von der Menge \(\Lambda_\sigma(F)\) der Labellings ab, nicht
von der Semantik \(\sigma\) selbst. Stimmen zwei Semantiken auf \(F\) in ihren
Labellingmengen \"uberein, so stimmen sie auch in s\"amtlichen
Unabh\"angigkeitsurteilen \"uberein.

\begin{Lemma}[Transferlemma]
    \label{lem:gcp-transfer}
    Sei \(F=(A,R)\) ein AF und seien \(\sigma,\tau\) zwei Semantiken mit
    \(\Lambda_\sigma(F)=\Lambda_\tau(F)\). Dann gilt f\"ur alle paarweise
    disjunkten Mengen \(X,Y,Z\subseteq A\):
    \[
        X\bot_\sigma Y\mid Z
        \quad\Longleftrightarrow\quad
        X\bot_\tau Y\mid Z.
    \]
\end{Lemma}

\begin{proof}
    Die definierende Bedingung von \(X\bot_\sigma Y\mid Z\) quantifiziert allein
    \"uber Elemente von \(\Lambda_\sigma(F)\): \enquote{f\"ur alle
    \(\lambda_1,\lambda_2\in\Lambda_\sigma(F)\) mit \(\lambda_1|_Z=\lambda_2|_Z\)
    existiert \(\lambda_3\in\Lambda_\sigma(F)\) mit \dots}. Ersetzt man
    \(\Lambda_\sigma(F)\) durch die nach Voraussetzung identische Menge
    \(\Lambda_\tau(F)\), so geht die Aussage Wort f\"ur Wort in die definierende
    Bedingung von \(X\bot_\tau Y\mid Z\) \"uber. Beide Seiten sind also
    \"aquivalent.
\end{proof}

Das Transferlemma ist das zentrale Werkzeug dieses Kapitels: Wann immer eine
Graphklasse \(\Lambda_{\mathrm{pr}}(F)=\Lambda_{\mathrm{st}}(F)\) erzwingt,
\"ubertragen sich \emph{s\"amtliche} Ergebnisse aus den
Kapiteln~\ref{ch:graphclasses} und~\ref{ch:graphclasses-empirical} unver\"andert
auf die \textit{preferred}-Semantik.

\subsection{Eindeutigkeit und fixierte Labels}
\label{subsec:gcp-eindeutig}

Das zweite Beweismittel verallgemeinert die f\"ur die \textit{stable}-Semantik
formulierte Proposition~\ref{prop:gce-leq1-indep} auf die
\textit{preferred}-Semantik.

\begin{Proposition}
    \label{prop:gcp-unique-indep}
    Sei \(F=(A,R)\) ein AF mit \(|\Lambda_{\mathrm{pr}}(F)|=1\). Dann gilt
    \(X\bot_{\mathrm{pr}}Y\mid Z\) f\"ur alle paarweise disjunkten Mengen
    \(X,Y,Z\subseteq A\).
\end{Proposition}

\begin{proof}
    Sei \(\Lambda_{\mathrm{pr}}(F)=\{\lambda\}\). F\"ur jedes Paar
    \(\lambda_1,\lambda_2\in\Lambda_{\mathrm{pr}}(F)\) gilt dann
    \(\lambda_1=\lambda_2=\lambda\). Mit \(\lambda_3:=\lambda\) sind
    \(\lambda_3|_X=\lambda_1|_X\), \(\lambda_3|_Y=\lambda_2|_Y\) und
    \(\lambda_3|_Z=\lambda_1|_Z\) unmittelbar erf\"ullt. Damit gilt
    \(X\bot_{\mathrm{pr}}Y\mid Z\).
\end{proof}

Anders als bei der \textit{stable}-Semantik ist hier
\(\Lambda_{\mathrm{pr}}(F)\neq\emptyset\) stets gesichert
\cite{dung:1995:af}. Die Bedingung \(|\Lambda_{\mathrm{pr}}(F)|\leq 1\) ist
daher gleichbedeutend mit \(|\Lambda_{\mathrm{pr}}(F)|=1\); einen vakuosen Fall
gibt es nicht. Dieser Unterschied wird in
Abschnitt~\ref{sec:gcp-scc} wieder aufgegriffen.

Oft gen\"ugt bereits, dass die fraglichen Argumente \"uber alle
\textit{preferred}-Labellings hinweg dasselbe Label tragen. Wir erinnern an das
eindeutig bestimmte \textit{grounded}-Labelling \(\lambda_{\mathrm{gr}}\) und
die bekannte Tatsache, dass es das informationsminimale
\textit{complete}-Labelling ist: F\"ur jedes
\(\lambda\in\Lambda_{\mathrm{co}}(F)\) gilt
\(\mathit{in}(\lambda_{\mathrm{gr}})\subseteq\mathit{in}(\lambda)\) und
\(\mathit{out}(\lambda_{\mathrm{gr}})\subseteq\mathit{out}(\lambda)\)
\cite{dung:1995:af,handbook__of_formal_argumentation__vol__3:2024}. Da
\(\Lambda_{\mathrm{pr}}(F)\subseteq\Lambda_{\mathrm{co}}(F)\), \"ubertr\"agt
sich dies auf alle \textit{preferred}-Labellings.

\begin{Proposition}
    \label{prop:gcp-fixed-label}
    Sei \(F=(A,R)\) ein AF und seien \(X,Y,Z\subseteq A\) paarweise disjunkt.
    Tr\"agt jedes Argument \(a\in X\) unter allen
    \(\lambda\in\Lambda_{\mathrm{pr}}(F)\) dasselbe Label, so gilt
    \(X\bot_{\mathrm{pr}}Y\mid Z\). Insbesondere gilt dies, falls
    \(X\subseteq \mathit{in}(\lambda_{\mathrm{gr}})\cup
    \mathit{out}(\lambda_{\mathrm{gr}})\) ist.
\end{Proposition}

\begin{proof}
    Seien \(\lambda_1,\lambda_2\in\Lambda_{\mathrm{pr}}(F)\) mit
    \(\lambda_1|_Z=\lambda_2|_Z\). W\"ahle \(\lambda_3:=\lambda_2\). Dann gilt
    \(\lambda_3|_Y=\lambda_2|_Y\) und
    \(\lambda_3|_Z=\lambda_2|_Z=\lambda_1|_Z\). Da jedes \(a\in X\) unter allen
    \textit{preferred}-Labellings dasselbe Label tr\"agt, gilt
    \(\lambda_2(a)=\lambda_1(a)\) f\"ur alle \(a\in X\), also
    \(\lambda_3|_X=\lambda_2|_X=\lambda_1|_X\). Somit ist \(\lambda_3\) ein
    Zeuge der Unabh\"angigkeit, und es gilt \(X\bot_{\mathrm{pr}}Y\mid Z\).

    Der Zusatz folgt, weil die Argumente in
    \(\mathit{in}(\lambda_{\mathrm{gr}})\) unter jedem
    \textit{preferred}-Labelling \(\mathit{in}\) und die Argumente in
    \(\mathit{out}(\lambda_{\mathrm{gr}})\) unter jedem
    \textit{preferred}-Labelling \(\mathit{out}\) sind.
\end{proof}

Proposition~\ref{prop:gcp-fixed-label} ist das mengenwertige Gegenst\"uck zu
Korollar~\ref{cor:singl-fixlabel} und erfasst insbesondere den
\enquote{grounded-Kern} eines AF, also den Teil, der bereits durch die
\textit{grounded}-Semantik entschieden ist.

\section{Azyklische AFs}
\label{sec:gcp-acyc}

\begin{Proposition}
    \label{prop:gcp-acyc-indep}
    Sei \(F=(A,R)\) ein azyklisches AF. Dann ist \(\textit{Ind}_{\mathrm{pr}}(F)\)
    stets wahr.
\end{Proposition}

\begin{proof}
    Ein endliches azyklisches AF ist wohlfundiert und besitzt daher genau ein
    \textit{complete}-Labelling, das zugleich \textit{grounded},
    \textit{preferred} und \textit{stable} ist
    \cite{dung:1995:af}. Folglich gilt \(|\Lambda_{\mathrm{pr}}(F)|=1\), und die
    Behauptung folgt mit Proposition~\ref{prop:gcp-unique-indep}.
\end{proof}

Diese Aussage ist das \textit{preferred}-Gegenst\"uck zu
Proposition~\ref{prop:acyc-af-ci-st-top}. Da das eindeutige Labelling mit dem
\textit{stable}- und dem \textit{grounded}-Labelling zusammenf\"allt, stimmen
auf azyklischen AFs s\"amtliche der drei Semantiken in ihren
Unabh\"angigkeitsurteilen \"uberein.

\section{Geradezykelfreie AFs}
\label{sec:gcp-even}

F\"ur die \textit{stable}-Semantik wurde in
Abschnitt~\ref{sec:even-cycle-free-afs} beobachtet, dass in AFs ohne gerade
Zyklen allein das \textit{grounded}-Labelling als \textit{stable}-Labelling in
Frage kommt. Strukturell liegt dem die st\"arkere Tatsache zugrunde, dass das
Vorhandensein eines geraden Zyklus notwendig ist, um \"uberhaupt mehrere
\textit{complete}-Labellings zu erzeugen; der kleinste Zeuge ist der
\(2\)-Zyklus \(a\leftrightarrow b\) mit den beiden Labellings
\(\mathit{in}(\lambda)=\{a\}\) bzw.\ \(\{b\}\).

\begin{Proposition}
    \label{prop:gcp-even-indep}
    Sei \(F=(A,R)\) ein AF ohne gerade Zyklen. Dann gilt
    \(|\Lambda_{\mathrm{co}}(F)|=1\); das eindeutige
    \textit{complete}-Labelling ist das \textit{grounded}-Labelling und zugleich
    das eindeutige \textit{preferred}-Labelling. Insbesondere ist
    \(\textit{Ind}_{\mathrm{pr}}(F)\) stets wahr.
\end{Proposition}

\begin{proof}
    Nach \cite{handbook__of_formal_argumentation__vol__3:2024} besitzt ein
    endliches AF ohne gerade Zyklen genau ein \textit{complete}-Labelling.
    Dieses ist als einziges zugleich \(\subseteq\)-minimal und
    \(\subseteq\)-maximal in der In-Menge, also \textit{grounded} und
    \textit{preferred}. Damit gilt \(|\Lambda_{\mathrm{pr}}(F)|=1\), und
    Proposition~\ref{prop:gcp-unique-indep} liefert die Unabh\"angigkeit.
\end{proof}

W\"ahrend f\"ur die \textit{stable}-Semantik in dieser Klasse \emph{h\"ochstens}
ein Labelling existiert (es kann auch keines geben), garantiert die
\textit{preferred}-Semantik \emph{genau} eines. Die Aussage ist damit das
sch\"arfere \textit{preferred}-Gegenst\"uck zu
Proposition~\ref{prop:even-cycle-free-af-ci-st-top}.

\section{Ungeradezykelfreie AFs}
\label{sec:gcp-odd}

F\"ur die \textit{stable}-Semantik bot diese Klasse keinen
berechnungstheoretischen Vorteil (Abschnitt~\ref{sec:odd-cycle-free-afs}). F\"ur
die \textit{preferred}-Semantik liegt der Fall g\"unstiger, weil das Fehlen
ungerader Zyklen die \textit{preferred}- und die \textit{stable}-Semantik zur
Deckung bringt.

\begin{Theorem}[Koh\"arenz, \citeauthor{dung:1995:af}]
    \label{th:gcp-coherent}
    Sei \(F=(A,R)\) ein endliches AF ohne ungerade Zyklen. Dann ist \(F\)
    koh\"arent, d.\,h.\ jedes \textit{preferred}-Labelling ist
    \textit{stable}, und es gilt
    \[
        \Lambda_{\mathrm{pr}}(F)=\Lambda_{\mathrm{st}}(F)\neq\emptyset.
    \]
\end{Theorem}

\begin{proof}[Beweisidee]
    Nach \textcite{dung:1995:af} ist jedes AF ohne ungerade Zyklen koh\"arent,
    es gilt also \(\Lambda_{\mathrm{pr}}(F)\subseteq\Lambda_{\mathrm{st}}(F)\).
    Umgekehrt ist jedes \textit{stable}-Labelling stets auch \textit{preferred}
    (eine \(\mathit{und}\)-freie In-Menge ist \(\subseteq\)-maximal), also
    \(\Lambda_{\mathrm{st}}(F)\subseteq\Lambda_{\mathrm{pr}}(F)\). Zusammen folgt
    Gleichheit. Die Nichtleerheit ergibt sich daraus, dass jedes AF ohne
    ungerade Zyklen mindestens ein \textit{stable}-Labelling besitzt
    \cite{dung:1995:af,handbook__of_formal_argumentation__vol__3:2024}.
\end{proof}

\begin{Corollary}
    \label{cor:gcp-odd-transfer}
    Sei \(F=(A,R)\) ein AF ohne ungerade Zyklen. Dann gilt f\"ur alle paarweise
    disjunkten \(X,Y,Z\subseteq A\):
    \[
        X\bot_{\mathrm{pr}}Y\mid Z
        \quad\Longleftrightarrow\quad
        X\bot_{\mathrm{st}}Y\mid Z.
    \]
    Insbesondere stimmen \(\textit{Ind}_{\mathrm{pr}}(F)\) und
    \(\textit{Ind}_{\mathrm{st}}(F)\) \"uberein.
\end{Corollary}

\begin{proof}
    Folgt aus Theorem~\ref{th:gcp-coherent} mit dem Transferlemma
    \ref{lem:gcp-transfer} f\"ur \(\sigma=\mathrm{pr}\) und
    \(\tau=\mathrm{st}\).
\end{proof}

Korollar~\ref{cor:gcp-odd-transfer} \"ubertr\"agt s\"amtliche Resultate der
\textit{stable}-Analyse auf die \textit{preferred}-Semantik. Komplexit\"atstheoretisch
sinkt damit die obere Schranke f\"ur \(\textit{Ind}_{\mathrm{pr}}\) auf dieser
Klasse von \(\Pi_3^{\mathrm P}\) auf \(\Pi_2^{\mathrm P}\): Die Frage l\"asst sich
\"aquivalent als \textit{stable}-Unabh\"angigkeit entscheiden
(Theorem~\ref{th:ind-st-is-pi2p-c}).

\section{Symmetrische AFs}
\label{sec:gcp-sym}

Auch symmetrische AFs f\"uhren auf eine Deckung von \textit{preferred}- und
\textit{stable}-Semantik, allerdings aus einem anderen Grund als im
ungeradezykelfreien Fall: Symmetrische AFs k\"onnen sehr wohl ungerade Zyklen
enthalten (etwa das symmetrische Dreieck), sind aber dennoch koh\"arent.

\begin{Theorem}[\citeauthor{handbook__of_formal_argumentation__vol__3:2024}]
    \label{th:gcp-sym}
    Sei \(F=(A,R)\) ein endliches, symmetrisches (\(R=R^{-1}\)) und
    irreflexives (\((a,a)\notin R\) f\"ur alle \(a\)) AF. Dann sind die
    zul\"assigen Mengen genau die konfliktfreien Mengen, und es gilt
    \[
        \Lambda_{\mathrm{pr}}(F)=\Lambda_{\mathrm{st}}(F),
    \]
    wobei die \textit{preferred}- und \textit{stable}-Labellings genau den
    \(\subseteq\)-maximalen konfliktfreien Mengen entsprechen.
\end{Theorem}

\begin{proof}[Beweisidee]
    In einem symmetrischen, irreflexiven AF verteidigt jede konfliktfreie Menge
    \(S\) sich selbst: Greift \(b\) ein \(a\in S\) an, so greift wegen der
    Symmetrie \(a\) auch \(b\) an, also wird \(a\) durch \(S\) verteidigt.
    Damit sind zul\"assige und konfliktfreie Mengen identisch, und die
    \(\subseteq\)-maximalen unter ihnen sind die \textit{preferred}-Extensionen.
    Eine maximale konfliktfreie Menge greift zudem jedes \"au\ss ere Argument an,
    ist also \textit{stable}. Daraus folgt
    \(\Lambda_{\mathrm{pr}}(F)=\Lambda_{\mathrm{st}}(F)\)
    \cite{handbook__of_formal_argumentation__vol__3:2024}.
\end{proof}

\begin{Corollary}
    \label{cor:gcp-sym-transfer}
    Sei \(F=(A,R)\) ein endliches, symmetrisches und irreflexives AF. Dann gilt
    f\"ur alle paarweise disjunkten \(X,Y,Z\subseteq A\):
    \[
        X\bot_{\mathrm{pr}}Y\mid Z
        \quad\Longleftrightarrow\quad
        X\bot_{\mathrm{st}}Y\mid Z.
    \]
\end{Corollary}

\begin{proof}
    Folgt aus Theorem~\ref{th:gcp-sym} mit dem Transferlemma
    \ref{lem:gcp-transfer}.
\end{proof}

Damit gelten alle Aussagen aus Abschnitt~\ref{sec:sym-afs} unver\"andert auch
f\"ur die \textit{preferred}-Semantik. Insbesondere \"ubertragen sich die in
Kapitel~\ref{ch:graphclasses-empirical} empirisch beobachteten Dichte-Muster
(Tabelle~\ref{tab:gce-sym}): Mit wachsender Kantenwahrscheinlichkeit \(p\)
steigt auch die \textit{preferred}-Abh\"angigkeitsrate, da die symmetrisch
erzeugten Instanzen irreflexiv sind und somit
\(\Lambda_{\mathrm{pr}}=\Lambda_{\mathrm{st}}\) erf\"ullen.

Als Extremfall halten wir den vollst\"andig gerichteten Graphen fest.

\begin{Corollary}
    \label{cor:gcp-cdg}
    Sei \(F=(A,R)\) ein vollst\"andig gerichtetes AF mit \(|A|\geq 2\), und
    seien \(X,Y,Z\subseteq A\) paarweise disjunkt mit \(X\neq\emptyset\) und
    \(Y\neq\emptyset\). Dann gilt \(X\not\bot_{\mathrm{pr}}Y\mid Z\).
\end{Corollary}

\begin{proof}
    Ein vollst\"andig gerichtetes AF ist symmetrisch und irreflexiv, also gilt
    nach Theorem~\ref{th:gcp-sym} \(\Lambda_{\mathrm{pr}}(F)=
    \Lambda_{\mathrm{st}}(F)\). Nach Lemma~\ref{lem:cdg-afs-labellings} sind dies
    genau die Labellings \(\lambda_a\) mit \(\mathit{in}(\lambda_a)=\{a\}\).
    Mit Proposition~\ref{prop:cdg-afs-not-st-ci} gilt
    \(X\not\bot_{\mathrm{st}}Y\mid Z\); das Transferlemma
    \ref{lem:gcp-transfer} liefert \(X\not\bot_{\mathrm{pr}}Y\mid Z\).
\end{proof}

\section{Stark verbundene Komponenten und der allgemeine Fall}
\label{sec:gcp-scc}

F\"ur die \textit{stable}-Semantik war die Klasse der dichten, stark
zusammenh\"angenden AFs dadurch gekennzeichnet, dass sie h\"aufig \emph{kein}
\textit{stable}-Labelling besitzt (Abschnitt~\ref{sec:scc-afs} sowie
Tabelle~\ref{tab:gce-scc}). Solche Instanzen waren \emph{vakuos} unabh\"angig:
Mangels Labellings ist die Unabh\"angigkeitsbedingung leer erf\"ullt.

F\"ur die \textit{preferred}-Semantik entf\"allt dieses Argument
vollst\"andig, denn \(\Lambda_{\mathrm{pr}}(F)\neq\emptyset\) ist stets
gesichert \cite{dung:1995:af}. Die folgende Aussage zeigt, dass die
vakuose Unabh\"angigkeit der \textit{stable}-Semantik unter der
\textit{preferred}-Semantik sogar in \emph{echte} Abh\"angigkeit umschlagen
kann.

\begin{Proposition}
    \label{prop:gcp-vacuous-vs-dep}
    Es gibt ein AF \(F=(A,R)\) mit \(\Lambda_{\mathrm{st}}(F)=\emptyset\) und
    Argumenten \(x,y\in A\), sodass
    \[
        \{x\}\bot_{\mathrm{st}}\{y\}\mid\emptyset
        \qquad\text{(vakuos)},
        \qquad\text{aber}\qquad
        \{x\}\not\bot_{\mathrm{pr}}\{y\}\mid\emptyset.
    \]
\end{Proposition}

\begin{proof}
    Wir betrachten das AF \(F=(A,R)\) mit
    \[
        A=\{a,b,c,d,e\},\qquad
        R=\{(a,b),(b,a)\}\cup\{(c,d),(d,e),(e,c)\}.
    \]
    Es besteht aus einem \(2\)-Zyklus \(a\leftrightarrow b\) und einem davon
    unabh\"angigen \(3\)-Zyklus \(c\to d\to e\to c\).

    \emph{Keine \textit{stable}-Labellings.} Ein \textit{stable}-Labelling
    h\"atte \(\mathit{und}=\emptyset\). Im \(3\)-Zyklus \(c\to d\to e\to c\)
    m\"usste dann jedes Argument \(\mathit{in}\) oder \(\mathit{out}\) sein; ein
    solches widerspruchsfreies Labelling existiert f\"ur einen ungeraden Zyklus
    jedoch nicht. Also gilt \(\Lambda_{\mathrm{st}}(F)=\emptyset\), und
    \(\{x\}\bot_{\mathrm{st}}\{y\}\mid\emptyset\) ist f\"ur beliebige \(x,y\)
    vakuos erf\"ullt.

    \emph{Genau zwei \textit{preferred}-Labellings.} Die Argumente
    \(c,d,e\) werden in jedem \textit{complete}-Labelling \(\mathit{und}\)
    gelabelt, da der ungerade Zyklus keine andere widerspruchsfreie Bewertung
    zul\"asst. Im \(2\)-Zyklus \(a\leftrightarrow b\) gibt es bei
    \(\subseteq\)-maximaler In-Menge genau die beiden Wahlen
    \(\mathit{in}(\lambda)=\{a\}\) oder \(\mathit{in}(\lambda)=\{b\}\). Somit
    ist \(\Lambda_{\mathrm{pr}}(F)=\{\lambda_a,\lambda_b\}\) mit
    \[
        \lambda_a:\ a\mapsto\mathit{in},\ b\mapsto\mathit{out},\quad
        \lambda_b:\ a\mapsto\mathit{out},\ b\mapsto\mathit{in},
    \]
    und \(c,d,e\mapsto\mathit{und}\) in beiden F\"allen.

    \emph{Abh\"angigkeit.} W\"ahle \(x:=a\) und \(y:=b\). Wegen
    \(Z=\emptyset\) gilt \(\lambda_a|_Z=\lambda_b|_Z\) trivial. Ein Zeuge
    \(\lambda_3\in\Lambda_{\mathrm{pr}}(F)\) m\"usste
    \(\lambda_3(a)=\lambda_a(a)=\mathit{in}\) und
    \(\lambda_3(b)=\lambda_b(b)=\mathit{in}\) erf\"ullen. Ein solches Labelling
    mit \(a,b\) gleichzeitig \(\mathit{in}\) existiert wegen \((a,b)\in R\)
    nicht. Also gilt \(\{a\}\not\bot_{\mathrm{pr}}\{b\}\mid\emptyset\).
\end{proof}

Proposition~\ref{prop:gcp-vacuous-vs-dep} belegt, dass sich die
\"ubergreifenden Muster der \textit{stable}-Analyse nicht naiv \"ubertragen
lassen. Insbesondere gilt das in Kapitel~\ref{ch:graphclasses-empirical}
beobachtete Muster \enquote{kaum Labellings, daher unabh\"angig} f\"ur die
\textit{preferred}-Semantik gerade \emph{nicht}: Wo die \textit{stable}-Semantik
mangels Labellings schweigt, kann die \textit{preferred}-Semantik eine echte
Abh\"angigkeit aufdecken.

F\"ur den allgemeinen Fall bleibt damit das Kriterium aus
Proposition~\ref{prop:gcp-unique-indep} bzw.\
Proposition~\ref{prop:gcp-fixed-label} ma\ss geblich: Generische
Unabh\"angigkeit liegt vor, wenn nur ein einziges
\textit{preferred}-Labelling existiert oder die fraglichen Argumente \"uber alle
\textit{preferred}-Labellings hinweg festgelegt sind. Au\ss erhalb der
strukturell verstandenen Klassen bleibt das Problem mit
\(\Pi_3^{\mathrm P}\)-Vollst\"andigkeit (Theorem~\ref{th:ind-pr-is-pi3p-c})
allerdings echt schwerer als im \textit{stable}-Fall.

\section{Zusammenfassung}
\label{sec:gcp-zusammenfassung}

Tabelle~\ref{tab:gcp-overview} fasst die beweisbaren Aussagen dieses Kapitels
zusammen. Zwei Beweismittel tragen die gesamte Analyse: das Transferlemma
\ref{lem:gcp-transfer}, das Klassen mit
\(\Lambda_{\mathrm{pr}}=\Lambda_{\mathrm{st}}\) auf die bereits verstandene
\textit{stable}-Analyse zur\"uckf\"uhrt, und das Eindeutigkeitskriterium
\ref{prop:gcp-unique-indep}, das Klassen mit einem einzigen
\textit{preferred}-Labelling abdeckt.

\begin{table}[htbp]
    \myfloatalign
    \begin{tabularx}{\textwidth}{lXl} \toprule
        Klasse & Struktur der \textit{preferred}-Labellings &
        Resultat \\ \midrule
        azyklisch & genau eines (\(=\mathrm{gr}=\mathrm{st}\)) &
            stets unabh\"angig \\
        geradezyklenfrei & genau eines (\(=\mathrm{gr}\)) &
            stets unabh\"angig \\
        ungeradezyklenfrei & \(\Lambda_{\mathrm{pr}}=\Lambda_{\mathrm{st}}
            \neq\emptyset\) & \(\textit{Ind}_{\mathrm{pr}}\equiv
            \textit{Ind}_{\mathrm{st}}\) \\
        symmetrisch (irreflexiv) & \(\Lambda_{\mathrm{pr}}=\Lambda_{\mathrm{st}}
            =\) max.\ konfliktfrei & \(\textit{Ind}_{\mathrm{pr}}\equiv
            \textit{Ind}_{\mathrm{st}}\) \\
        vollst\"andig gerichtet & Singletons \(\{a\}\) &
            abh\"angig (\(X,Y\neq\emptyset\)) \\
        SCC / allgemein & stets \(\neq\emptyset\); keine Vakuit\"at &
            Kriterium \(|\Lambda_{\mathrm{pr}}|=1\) \\
        \bottomrule
    \end{tabularx}
    \caption[Graphklassen unter \textit{preferred}-Unabh\"angigkeit]{Beweisbare
    Aussagen zur \textit{preferred}-Unabh\"angigkeit nach Graphklassen.}
    \label{tab:gcp-overview}
\end{table}

Die wichtigste qualitative Erkenntnis betrifft die stark zusammenh\"angenden
Komponenten: W\"ahrend die \textit{stable}-Semantik dort h\"aufig vakuos
unabh\"angig war, sichert die garantierte Existenz von
\textit{preferred}-Labellings, dass dieselben Instanzen echte Abh\"angigkeiten
aufweisen k\"onnen (Proposition~\ref{prop:gcp-vacuous-vs-dep}). Komplexit\"atstheoretisch
spiegelt sich dies in der h\"oheren Einordnung von
\(\textit{Ind}_{\mathrm{pr}}\) (\(\Pi_3^{\mathrm P}\)) wider, die nur auf den
strukturell beherrschten Klassen --- azyklisch und geradezyklenfrei
(trivial entscheidbar) sowie ungeradezyklenfrei und symmetrisch (Reduktion auf
\(\Pi_2^{\mathrm P}\)) --- nachweislich sinkt.
